%! AP Statistics — Unit 8, Lesson 8.2 — Guided Notes (Answer Key)
\documentclass[10pt]{article}
\usepackage{apstats-article}
\usepackage{apstats-key}

\begin{document}

\pageheader{Unit 8, Lesson 8.2}{Guided Notes --- Answer Key}
\namedateperiod

\begin{objectivebox}
By the end of this lesson, I will be able to:
\begin{itemize}
    \item[$\square$] \ansline{Describe the shape of a chi-square distribution and the effect of degrees of freedom.}
    \item[$\square$] \ansline{State H$_0$ and H$_a$ for a chi-square goodness of fit test.}
    \item[$\square$] \ansline{Identify when a chi-square GOF test is the appropriate procedure.}
    \item[$\square$] \ansline{Calculate expected counts using $E = n \cdot p_0$.}
    \item[$\square$] \ansline{Verify the Random and Large Counts conditions for a GOF test.}
\end{itemize}
\end{objectivebox}

\begin{vocabbox}
\termblanklong{Chi-square distribution}
\ansline{A family of distributions with only positive values, skewed right; becomes less skewed as degrees of freedom increase.}

\termblanklong{Goodness of fit (GOF) test}
\ansline{A chi-square procedure that tests whether observed data for one categorical variable follow a claimed (specified) distribution.}

\termblanklong{Expected count}
\ansline{$E = n \cdot p_0$; the count expected in a category if H$_0$ is true.}

\termblanklong{Null proportion ($p_0$)}
\ansline{The claimed proportion for each category, specified by H$_0$.}

\termblanklong{Degrees of freedom ($df$)}
\ansline{For a GOF test: $df = (\text{number of categories}) - 1$.}

\termblanklong{Large Counts condition}
\ansline{All expected counts must be $\ge 5$.}
\end{vocabbox}

\begin{hookbox}
Why can't we use six separate one-proportion $z$-tests to check the die distribution?
\ansline{Using multiple separate tests inflates the overall Type I error rate. We need a single procedure that combines all six categories into one test statistic --- the chi-square GOF test.}
\end{hookbox}

\begin{notesbox}{The Chi-Square Distribution (VAR-8.A)}
\textbf{Properties:}
\begin{itemize}
    \item Values are always \ans{positive (or zero)}.
    \item Shape is skewed \ans{right}.
    \item As degrees of freedom \ans{increase}, the skew becomes less pronounced.
    \item $df = (\text{number of categories}) - $ \ans{1} for a GOF test.
\end{itemize}

\textbf{When to use a GOF test (VAR-8.C):} Use a chi-square GOF test when you have
\ans{one} categorical variable and want to test whether the population follows a
\ans{specified (claimed)} distribution.
\end{notesbox}

\begin{notesbox}{Stating Hypotheses (VAR-8.B)}
H$_0$ states the \ans{proportion (null proportion)} for \ans{each} category.

H$_a$: At least \ans{one} of the proportions is not as specified in H$_0$.

\vspace{0.1in}
\textbf{Die-Roll Example:}

H$_0$: $p_1 = $ \ans{$1/6$}, $p_2 = $ \ans{$1/6$}, $p_3 = $ \ans{$1/6$},
$p_4 = $ \ans{$1/6$}, $p_5 = $ \ans{$1/6$}, $p_6 = $ \ans{$1/6$}

H$_a$: \ansline{At least one face of the die does not appear with probability $1/6$.}

\textit{Note:} All null proportions must sum to \ans{1}.
\end{notesbox}

\begin{notesbox}{Calculating Expected Counts (VAR-8.D)}
Formula: \quad $E = $ \ans{sample size ($n$)} $\times$ \ans{null proportion ($p_0$)}

\vspace{0.1in}
\textbf{Die-Roll Example} ($n = 120$): \quad $E = 120 \times \dfrac{1}{6} = $ \ans{20} per face.

\vspace{0.1in}
\begin{center}
\renewcommand{\arraystretch}{1.4}
\begin{tabular}{lcccccc}
    \textbf{Face} & 1 & 2 & 3 & 4 & 5 & 6 \\
    \hline
    \textbf{Null proportion} & $1/6$ & $1/6$ & $1/6$ & $1/6$ & $1/6$ & $1/6$ \\
    \textbf{Expected count} & \ans{20} & \ans{20} & \ans{20} & \ans{20} & \ans{20} & \ans{20} \\
\end{tabular}
\end{center}
\end{notesbox}

\begin{notesbox}{Verifying Conditions (VAR-8.E)}
\begin{enumerate}
    \item \textbf{Random:} The data come from a \ans{random sample} or randomized experiment.
          If sampling without replacement: $n \le $ \ans{10}\%$ \cdot N$.
    \item \textbf{Large Counts:} All \ans{expected} counts are $\ge $ \ans{5}.
\end{enumerate}

\textbf{Die-Roll Check:}
\begin{itemize}
    \item Random: \ansline{The 120 rolls are assumed to be a random process (fair rolling conditions).}
    \item Large Counts: Smallest expected count is \ans{20}. Condition \ans{met ($\ge 5$)}.
\end{itemize}

\textit{Important:} The Large Counts condition uses \ans{expected} counts, not
\ans{observed} counts.
\end{notesbox}

\begin{practicebox}
\textbf{Practice:} Total ratio parts: $9 + 3 + 3 + 1 = $ \ans{16}

\begin{center}
\renewcommand{\arraystretch}{1.4}
\begin{tabular}{lcccc}
    \textbf{Phenotype} & A & B & C & D \\
    \hline
    \textbf{Null proportion} & $9/16$ & $3/16$ & $3/16$ & $1/16$ \\
    \textbf{Expected count} ($n=160$) & \ans{90} & \ans{30} & \ans{30} & \ans{10} \\
\end{tabular}
\end{center}

Large Counts condition: \ansline{All expected counts $\ge 5$ (smallest is 10). Condition is met.}
\end{practicebox}

\end{document}
